\documentclass[11pt,a4paper,sans]{moderncv}

% moderncv themes
\moderncvstyle{classic}
\moderncvcolor{blue}

% character encoding
\usepackage[utf8]{inputenc}
\usepackage{xeCJK}

% adjust the page margins
\usepackage[scale=0.88]{geometry}

% personal data
\name{}{庞朝阳}
\title{AI 应用开发工程师}
\address{广西外国语学院}{计算机科学与技术}{2026 届本科}
\phone[mobile]{18977554553}
\email{2380691181@qq.com}
\social[github]{pofice}

%----------------------------------------------------------------------------------
%            content
%----------------------------------------------------------------------------------
\begin{document}
\makecvtitle

\section{个人简介}
\cvitem{}{2026 届 CS 本科毕业，专注 AI Native 全栈工程——不只使用 AI 工具，更构建 AI Agent 可调用的工具、运行时与决策规范。现就职于雷石（KTV 行业龙头）AI 团队（小米总部驻场），主力开发 AI 客服 Agent 平台与 ERP 多端产品矩阵，另主导过跨境短剧 AIGC 系统技术架构。}

\section{教育背景}
\cventry{2022--2026}{本科}{广西外国语学院}{计算机科学与技术}{}{}

\section{专业技能}
\cvitem{主力}{TypeScript / Node.js, Python, Vue 3}
\cvitem{熟悉}{Go (Gin / GORM / DDD 四层), FastAPI, uni-app}
\cvitem{AI 工程}{Claude Code（自研 Skills 与 CLAUDE.md 规范）, MCP, 多 Agent SDK 集成与 Agent Runtime, RAG, Prompt 工程, Dify / Coze / n8n}

\section{工作经历}
\cventry{2026.01--至今}{AI 应用开发工程师（小米总部驻场）}{北京雷石天地电子技术股份有限公司 / AI 团队}{北京}{}{
\begin{itemize}
    \item 主力开发企业级 AI 客服 Agent 平台「小雷」（TypeScript / Node.js + Fastify）：微信 + 飞书双渠道消息网关（RabbitMQ + WebSocket 长连接），Orchestrator 调度可插拔多 Agent 后端（Claude Code / Kimi / OpenCode SDK），通过 MCP 挂载业务工具，配套知识库系统（100+ 篇文档）与按群聊作用域的知识隔离；已投产承接 VOD 产品线技术支持——持续运行 4 个多月，覆盖 24 个业务群 / 300+ 用户，累计处理咨询消息 2.8 万+ 条、自动应答近 5,000 条（近 30 天日均入站 400+）；8,700+ 次应答决策中 95\% 判定"无需介入"——把群聊 Agent 的克制做成产品能力，全链路审计留痕（4.6 万+ 审计事件）。
    \item 从 0 到 1 交付 VOD 自动退款 Agent（微信消息采集 → OCR 提取订单号 → CMS 接口对接 → 飞书告警），已投产替代人工退款流程；工程化踩坑沉淀：并发锁防双倍退款、已退款订单幂等判断、异常容错重试。
    \item 深度参与 KTV ERP 多端产品矩阵开发——客用小程序、掌柜端小程序（uni-app + Vue3）为主力，兼及收银端 / 包厢 Pad、财务报表系统、SaaS 管理后台（Vue3）与 Go 后端，交付自助预约、扫码开台、卡券核销、超时自动关台等完整业务闭环；攻克跨端数据一致性问题（预约时段冲突、session 状态漂移），抽象为状态机模型在领域服务层统一修复。
    \item 深入 Go DDD 四层架构后端（规则引擎 / 流程引擎），承担订单、预约、会员等核心交易链路的常规迭代（270+ commits），配套单元测试与集成测试。
    \item 为客用端 / 掌柜端建立 Playwright e2e 自动化验收体系（40+ 场景、60+ 用例）：真实 H5 UI 驱动 + MySQL 状态断言，覆盖预约、开台、换房、权限等核心链路，作为发版闸门替代人工回归；日常开发全面接入 Claude Code，以 CLAUDE.md 规范沉淀代码库上下文与团队踩坑经验。
\end{itemize}}

\cventry{2025.07--2025.09}{全栈工程师（实习）}{北京顺为知一科技有限公司}{北京}{}{
\begin{itemize}
    \item 负责企业级微信销售助手平台全栈开发（FastAPI + Next.js），支持多微信账号统一接入、自动化消息处理与客户标签体系；通信层基于开源 wxauto 库扩展开发。
    \item 将 AI 能力（智能对话、内容摘要、自动标签）下沉到 Dify / Coze / n8n 工作流，后端仅保留调度与审计，使运营侧可自主迭代 Prompt 而无需后端发版。
    \item RabbitMQ + WebSocket 构建实时消息推送链路，支撑多账号并发场景。
\end{itemize}}

\newpage
\section{开源项目}
\cventry{2024.09--至今}{voice-input-method — Linux 离线语音输入法}{\textnormal{GitHub \textbf{29} stars \textperiodcentered{} 9 forks}}{Python}{}{
\begin{itemize}
    \item 基于 FunASR / SherpaOnnx + PySide6 的跨平台离线语音输入工具，响应延迟 \textless{} 1s，支持 X11 / Wayland 双协议与运行时麦克风切换；核心方案（常驻音频流消除录音启动延迟 + 内存/磁盘双写 + 二遍识别）已申请发明专利。
    \item Protocol-based 依赖注入架构（GUI → factory → engine → protocols → implementations），核心逻辑不依赖硬件即可 mock 测试。配套 CLAUDE.md 规范使 AI Agent 可直接接管代码库——新贡献者（人或 AI）无需口头说明即可上手开发。
\end{itemize}}

\cventry{2026.02--至今}{mc-intelligence — 让 AI Agent 在 Minecraft 里实时建造}{\textnormal{NeoForge mod + CLI + MCP（前身 minecraft-builder Skill, GitHub 7 stars）}}{Java / Python}{}{
\begin{itemize}
    \item 设计理念：不是"用 AI 写代码"，而是"为 AI 设计可安全调用的工具"。从离线存档工具演进为实时运行环境：NeoForge mod 内嵌 HTTP 服务器暴露建造原语，配套 CLI 与 MCP 接口，AI Agent 可在游戏运行中零人工干预完成地形 / 建筑 / 家具生成；将 LLM 高频幻觉点（块名拼写、坐标越界、存档损坏）内置到防御层，配套 CLAUDE.md + SKILL.md 规范与 CI 覆盖率门禁。
\end{itemize}}

\section{项目经历}
\cventry{2026.03--至今}{观澜 — 黄金交易分析 Claude Code Skill}{\textnormal{412 行结构化规范 \textperiodcentered{} AI Agent 决策工程化}}{}{}{
\begin{itemize}
    \item 核心命题：如何让 LLM 在高风险领域（金融）输出可信结论？设计四层递进分析框架（周期定阶段 → 宏观定方向 → 中观定节奏 → 微观定点位），将交易专家认知系统化为 AI Agent 可执行的决策规范，附优先级铁律：上层逻辑 > 下层信号，矛盾时必须解释数据依据。
    \item 工程化约束 AI 幻觉：制定"数据锚点原则"——每个结论必须挂钩数据源 + 具体值 + 时间戳（三位一体），禁止无数据支撑的判断性形容词；设计八因子量化评分表（$-$16 到 $+$16）集成 FRED / yfinance / CFTC COT 等多源数据，并构建中观六维信号交叉验证矩阵（COMEX 持仓 + ETF 流向 + 央行购金 + 交割博弈 + 波动率 + 期权），用组合模式替代单一信号判断。
\end{itemize}}

\cventry{2023--2024}{网易我的世界高铁模组}{\textnormal{累计下载 \textbf{10,000+}}}{python}{}{
\begin{itemize}
    \item 与同学组队开发网易我的世界高铁模组，负责核心功能开发与模组架构设计；发布至网易我的世界模组平台后持续维护迭代，累计下载量突破 10,000 次，累计营收突破 100,000 元。
\end{itemize}}

\cventry{2025.09--2026.01}{技术负责人 — 跨境短剧 AIGC 系统}{映乾（3 人创业团队，负责全部技术架构与 AI 工作流）}{}{}{
\begin{itemize}
    \item 主导技术选型与系统架构设计：Sora2 API + Dify 流式剧本编排 + Celery + Redis 异步任务链路（FastAPI + Vue3 + PostgreSQL + Docker）；设计角色镜头截取机制，解决多段 AI 生成视频的角色一致性问题。
    \item 独立搭建翻译配音全自动链路：Demucs v4 人声分离 → VieNeu-TTS / Index-TTS 音色克隆 → OCR / FunASR 双路字幕识别，实现上传到成片全流程无人工干预，单集处理耗时从人工 4h+ 降至自动 20min。
\end{itemize}}

\section{专利、竞赛与荣誉}
\cvitem{2026}{发明专利（独立发明人，已受理）：一种基于常驻音频流的低延迟语音输入方法及系统}
\cvitem{2025.10}{第十届中国国际大学生创新大赛 区级铜奖（队长）}
\cvitem{2025.06}{第十二届"挑战杯"广西大学生课外学术科技作品竞赛 一等奖（队长）}
\cvitem{2023.06}{大学生创新创业项目 国家级立项（负责人）}

\end{document}
